\documentclass[a4paper,11pt]{article}
\usepackage[utf8]{inputenc}
\usepackage[T1]{fontenc}
\usepackage[french]{babel}
\usepackage{geometry}
\geometry{left=1.5cm,right=1.5cm,top=1.2cm,bottom=1.2cm}

\usepackage{lmodern}
\usepackage{xcolor}
\definecolor{titleblue}{RGB}{0,51,140}

\usepackage{hyperref}
\hypersetup{colorlinks=true, urlcolor=titleblue, linkcolor=titleblue}

\usepackage{titlesec}
\usepackage{enumitem}

\pagestyle{empty}

\titleformat{\section}{\color{titleblue}\Large\bfseries\uppercase}{}{0em}{}[\color{titleblue}\titlerule]
\titlespacing\section{0pt}{8pt}{4pt}
\setlist[itemize]{leftmargin=*,topsep=1pt,itemsep=1pt}

\begin{document}

\begin{center}
    {\Huge\bfseries\color{titleblue} AISSAOUI Ismail}\\[4pt]
    {\Large Ingénieur d'État en Intelligence Artificielle \& Data Science}\\[6pt]
    ismail.aissaoui.pro@gmail.com $|$ +213 660 70 77 96 $|$ Algérie \\
    \href{https://linkedin.com/in/aissaoui-ismail-6a77a92a7}{linkedin.com/in/aissaoui-ismail-6a77a92a7} $|$ \href{https://ismail-aissaoui.vercel.app}{ismail-aissaoui.vercel.app}
\end{center}

\section{RÉSUMÉ DE RECHERCHE}
Ingénieur d'État (ESI-SBA) spécialisé dans les \textbf{Jumeaux Numériques}, l'\textbf{Exploration d'Espaces de Conception (DSE)} et l'\textbf{Optimisation Multi-Objectif}. Expert dans la résolution de compromis non-fonctionnels (Énergie, Latence, Précision) pour les systèmes d'IA industriels. Maîtrise des algorithmes métaheuristiques (PSO, GA) pour la configuration automatisée et la synchronisation de systèmes distribués.

\section{FORMATION}
\textbf{École Nationale Supérieure d'Informatique (ESI-SBA, Algérie)} \hfill 2021 -- 2026 \\
\textbf{Diplôme d'Ingénieur d'État \& Master en Intelligence Artificielle}

\section{PROJETS CLÉS}

\textbf{Jumeau Numérique co-conçu Matériel-Logiciel — Sonatrach} \hfill 2026 \\
Optimisation de substituts génératifs pour environnements industriels contraints.
\begin{itemize}
    \item Développement d'une architecture hybride \textbf{Mamba-SSM/xLSTM} optimisée pour les CPU Edge.
    \item Résolution du paradoxe latence-précision avec une inférence de \textbf{0.04ms} via compilation \textbf{C++ JIT}.
    \item Implémentation de contraintes \textbf{PINN} pour garantir la cohérence physique des simulations.
\end{itemize}

\textbf{NLP Avancé : Optimisation Métaheuristique de Modèles} \hfill 2025 \\
Plateforme de réglage automatisé des hyperparamètres et d'exploration d'architectures.
\begin{itemize}
    \item Implémentation de l'\textbf{Optimisation par Essaim Particulaire (PSO)} et d'\textbf{Algorithmes Génétiques (GA)} pour explorer l'espace de recherche de DistilBERT.
    \item Développement de visualisations en temps réel pour la convergence et l'exploration du front de Pareto.
\end{itemize}

\textbf{Système POS Hybride \& Moteur de Sync E-commerce} \hfill 2025 \\
Système distribué de classe entreprise avec réconciliation de données en temps réel.
\begin{itemize}
    \item Architecture d'une **couche double base de données** (SQLite/PostgreSQL) pour la résilience hors-ligne.
    \item Développement de microservices de synchronisation gérant la cohérence des stocks multi-boutiques.
\end{itemize}

\section{COMPÉTENCES TECHNIQUES}
\begin{itemize}
    \item \textbf{Architecture} : Pipelines de Jumeaux Numériques, Systèmes Distribués, Electron, Next.js.
    \item \textbf{Optimisation} : PSO, Algorithmes Génétiques, Optimisation Bayésienne, Optuna.
    \item \textbf{Génie Logiciel ML} : PyTorch, C++ JIT, Docker, PostgreSQL, SQLite, FAISS.
    \item \textbf{Systèmes de Données} : Synchronisation en temps réel, Bases de données vectorielles (Chroma).
\end{itemize}

\section{LANGUES}
Arabe (Maternel), Anglais (Avancé - Recherche Professionnelle), Français (Courant).

\end{document}
